\chapter{Complete prompt texts}

The following are the nine exact main-condition prompts for example\_\allowbreak{}0984, framboid, as stored in the supplied Gemini records. They include the external demonstration, JSON response instructions, and graph notes. Text is wrapped for the page; wording and punctuation are retained. In these archived prompts, “target” refers to the omitted centre class; it does not redefine Gmax(e). This makes graph/no-graph template differences directly inspectable.


\FloatBarrier

\section{Standard no graph}

\begin{Verbatim}[fontsize=\footnotesize,breaklines=true,breakanywhere=true]
List valid entity-type labels, most specific first, then broader valid types. Use concise labels, not entity names or descriptions. Include only evidence-supported types; do not repeat labels or pad the list.
\end{Verbatim}

\begin{Verbatim}[fontsize=\footnotesize,breaklines=true,breakanywhere=true]
Return JSON only in exactly this form:
{"types": ["type 1", "type 2", "..."]}
\end{Verbatim}

\begin{Verbatim}[fontsize=\footnotesize,breaklines=true,breakanywhere=true]
Example:
\end{Verbatim}

\begin{Verbatim}[fontsize=\footnotesize,breaklines=true,breakanywhere=true]
Entity: Rod Fergusson
Description: Canadian video game producer
\end{Verbatim}

\begin{Verbatim}[fontsize=\footnotesize,breaklines=true,breakanywhere=true]
Answer:
{"types": ["video game producer", "producer", "media professional", "professional", "worker", "person", "person or organization", "corporate body", "object", "entity"]}
\end{Verbatim}

\begin{Verbatim}[fontsize=\footnotesize,breaklines=true,breakanywhere=true]
Now classify:
\end{Verbatim}

\begin{Verbatim}[fontsize=\footnotesize,breaklines=true,breakanywhere=true]
Entity: framboid
Description: minute spheroidal grain of pyrite or an aggregate of such grains, commonly formed at low temperatures in oxygen-poor mud sediments containing organic matter, especially in marine mud deposits
\end{Verbatim}

\begin{Verbatim}[fontsize=\footnotesize,breaklines=true,breakanywhere=true]
Answer:
\end{Verbatim}

\clearpage

\section{Standard graph}

\begin{Verbatim}[fontsize=\footnotesize,breaklines=true,breakanywhere=true]
List valid entity-type labels, most specific first, then broader valid types. Use concise labels, not entity names or descriptions. Include only evidence-supported types; do not repeat labels or pad the list.
\end{Verbatim}

\begin{Verbatim}[fontsize=\footnotesize,breaklines=true,breakanywhere=true]
The context surrounds an omitted target type; it is evidence, not answer options. Parents are broader, children narrower, and siblings share a parent. Infer the type using the description and these relations. Copy a neighbour only if valid for this entity. Siblings need not be disjoint; predict a child only when the description independently supports it.
\end{Verbatim}

\begin{Verbatim}[fontsize=\footnotesize,breaklines=true,breakanywhere=true]
Return JSON only in exactly this form:
{"types": ["type 1", "type 2", "..."]}
\end{Verbatim}

\begin{Verbatim}[fontsize=\footnotesize,breaklines=true,breakanywhere=true]
Example:
\end{Verbatim}

\begin{Verbatim}[fontsize=\footnotesize,breaklines=true,breakanywhere=true]
Entity: Rod Fergusson
Description: Canadian video game producer
\end{Verbatim}

\begin{Verbatim}[fontsize=\footnotesize,breaklines=true,breakanywhere=true]
Taxonomy context:
- Direct parent classes: producer
- Sibling classes: creative producer, co-producer, line producer
- Direct child classes: None
\end{Verbatim}

\begin{Verbatim}[fontsize=\footnotesize,breaklines=true,breakanywhere=true]
Answer:
{"types": ["video game producer", "producer", "media professional", "professional", "worker", "person", "person or organization", "corporate body", "object", "entity"]}
\end{Verbatim}

\begin{Verbatim}[fontsize=\footnotesize,breaklines=true,breakanywhere=true]
Now classify:
\end{Verbatim}

\begin{Verbatim}[fontsize=\footnotesize,breaklines=true,breakanywhere=true]
Entity: framboid
Description: minute spheroidal grain of pyrite or an aggregate of such grains, commonly formed at low temperatures in oxygen-poor mud sediments containing organic matter, especially in marine mud deposits
\end{Verbatim}

\begin{Verbatim}[fontsize=\footnotesize,breaklines=true,breakanywhere=true]
Taxonomy context:
- Direct parent classes: chemical substance
- Sibling classes: chemical product, pesticide, aroma compound
- Direct child classes: None
\end{Verbatim}

\begin{Verbatim}[fontsize=\footnotesize,breaklines=true,breakanywhere=true]
Answer:
\end{Verbatim}

\clearpage

\section{CoT no graph}

\begin{Verbatim}[fontsize=\footnotesize,breaklines=true,breakanywhere=true]
List valid entity-type labels, most specific first, then broader valid types. Use concise labels, not entity names or descriptions. Include only evidence-supported types; do not repeat labels or pad the list.
\end{Verbatim}

\begin{Verbatim}[fontsize=\footnotesize,breaklines=true,breakanywhere=true]
Briefly explain the most specific supported type, reject unsupported nearby alternatives, and justify the broader types and their order. Use only the entity and description; do not assume supplied taxonomy evidence.
\end{Verbatim}

\begin{Verbatim}[fontsize=\footnotesize,breaklines=true,breakanywhere=true]
Return JSON only in exactly this form:
{"explanation": "description-grounded explanation", "types": ["type 1", "type 2", "..."]}
\end{Verbatim}

\begin{Verbatim}[fontsize=\footnotesize,breaklines=true,breakanywhere=true]
Example:
\end{Verbatim}

\begin{Verbatim}[fontsize=\footnotesize,breaklines=true,breakanywhere=true]
Entity: Rod Fergusson
Description: Canadian video game producer
\end{Verbatim}

\begin{Verbatim}[fontsize=\footnotesize,breaklines=true,breakanywhere=true]
Correct response:
{"explanation": "The description \"Canadian video game producer\" directly supports video game producer as the most specific type. Plausible alternatives that require details absent from the description should be rejected. The remaining returned types are broader categories that still apply, so the ranking moves from specific to broad.", "types": ["video game producer", "producer", "media professional", "professional", "worker", "person", "person or organization", "corporate body", "object", "entity"]}
\end{Verbatim}

\begin{Verbatim}[fontsize=\footnotesize,breaklines=true,breakanywhere=true]
Now classify:
\end{Verbatim}

\begin{Verbatim}[fontsize=\footnotesize,breaklines=true,breakanywhere=true]
Entity: framboid
Description: minute spheroidal grain of pyrite or an aggregate of such grains, commonly formed at low temperatures in oxygen-poor mud sediments containing organic matter, especially in marine mud deposits
\end{Verbatim}

\begin{Verbatim}[fontsize=\footnotesize,breaklines=true,breakanywhere=true]
Answer:
\end{Verbatim}

\clearpage

\section{CoT graph}

\begin{Verbatim}[fontsize=\footnotesize,breaklines=true,breakanywhere=true]
List valid entity-type labels, most specific first, then broader valid types. Use concise labels, not entity names or descriptions. Include only evidence-supported types; do not repeat labels or pad the list.
\end{Verbatim}

\begin{Verbatim}[fontsize=\footnotesize,breaklines=true,breakanywhere=true]
The context surrounds an omitted target type; it is evidence, not answer options. Parents are broader, children narrower, and siblings share a parent. Infer the type using the description and these relations. Copy a neighbour only if valid for this entity. Siblings need not be disjoint; predict a child only when the description independently supports it.
\end{Verbatim}

\begin{Verbatim}[fontsize=\footnotesize,breaklines=true,breakanywhere=true]
Briefly use parents to explain scope, compare sibling alternatives, and check children for unsupported specificity. Justify the final specific-to-broad order using the description and taxonomy evidence.
\end{Verbatim}

\begin{Verbatim}[fontsize=\footnotesize,breaklines=true,breakanywhere=true]
Return JSON only in exactly this form:
{"explanation": "taxonomy-grounded explanation", "types": ["type 1", "type 2", "..."]}
\end{Verbatim}

\begin{Verbatim}[fontsize=\footnotesize,breaklines=true,breakanywhere=true]
Example:
\end{Verbatim}

\begin{Verbatim}[fontsize=\footnotesize,breaklines=true,breakanywhere=true]
Entity: Rod Fergusson
Description: Canadian video game producer
\end{Verbatim}

\begin{Verbatim}[fontsize=\footnotesize,breaklines=true,breakanywhere=true]
Taxonomy context:
- Direct parent classes: producer
- Sibling classes: creative producer, co-producer, line producer
- Direct child classes: None
\end{Verbatim}

\begin{Verbatim}[fontsize=\footnotesize,breaklines=true,breakanywhere=true]
Correct response:
{"explanation": "The description \"Canadian video game producer\" supports video game producer as the most specific type. In the taxonomy, its direct parent is producer, which constrains the broader semantic category. The sibling classes are nearby alternatives, but classes such as creative producer, co-producer, line producer are not as well supported by the description. The listed child context is None; a child should be selected only if the description independently supports that narrower meaning. The remaining returned types are broader valid classes ranked in decreasing specificity.", "types": ["video game producer", "producer", "media professional", "professional", "worker", "person", "person or organization", "corporate body", "object", "entity"]}
\end{Verbatim}

\begin{Verbatim}[fontsize=\footnotesize,breaklines=true,breakanywhere=true]
Now classify:
\end{Verbatim}

\begin{Verbatim}[fontsize=\footnotesize,breaklines=true,breakanywhere=true]
Entity: framboid
Description: minute spheroidal grain of pyrite or an aggregate of such grains, commonly formed at low temperatures in oxygen-poor mud sediments containing organic matter, especially in marine mud deposits
\end{Verbatim}

\begin{Verbatim}[fontsize=\footnotesize,breaklines=true,breakanywhere=true]
Taxonomy context:
- Direct parent classes: chemical substance
- Sibling classes: chemical product, pesticide, aroma compound
- Direct child classes: None
\end{Verbatim}

\begin{Verbatim}[fontsize=\footnotesize,breaklines=true,breakanywhere=true]
Answer:
\end{Verbatim}

\clearpage

\section{Contrastive no graph}

\begin{Verbatim}[fontsize=\footnotesize,breaklines=true,breakanywhere=true]
List valid entity-type labels, most specific first, then broader valid types. Use concise labels, not entity names or descriptions. Include only evidence-supported types; do not repeat labels or pad the list.
\end{Verbatim}

\begin{Verbatim}[fontsize=\footnotesize,breaklines=true,breakanywhere=true]
Briefly explain the most specific supported type, reject unsupported nearby alternatives, and justify the broader types and their order. Use only the entity and description; do not assume supplied taxonomy evidence.
\end{Verbatim}

\begin{Verbatim}[fontsize=\footnotesize,breaklines=true,breakanywhere=true]
Contrast the best-supported type with plausible incorrect alternatives. Explain why unsupported, overly narrow, or incompatible alternatives are rejected.
\end{Verbatim}

\begin{Verbatim}[fontsize=\footnotesize,breaklines=true,breakanywhere=true]
Return JSON only in exactly this form:
{"explanation": "contrastive description-grounded explanation", "types": ["type 1", "type 2", "..."]}
\end{Verbatim}

\begin{Verbatim}[fontsize=\footnotesize,breaklines=true,breakanywhere=true]
Example:
\end{Verbatim}

\begin{Verbatim}[fontsize=\footnotesize,breaklines=true,breakanywhere=true]
Entity: Rod Fergusson
Description: Canadian video game producer
\end{Verbatim}

\begin{Verbatim}[fontsize=\footnotesize,breaklines=true,breakanywhere=true]
Correct response:
{"explanation": "The description \"Canadian video game producer\" directly supports video game producer as the most specific type. Plausible alternatives that require details absent from the description should be rejected. The remaining returned types are broader categories that still apply, so the ranking moves from specific to broad.", "types": ["video game producer", "producer", "media professional", "professional", "worker", "person", "person or organization", "corporate body", "object", "entity"]}
\end{Verbatim}

\begin{Verbatim}[fontsize=\footnotesize,breaklines=true,breakanywhere=true]
Wrong ranking:
{"types": ["producer", "creative producer", "co-producer", "line producer", "radio producer", "DJ producer", "professional", "worker", "person", "entity"]}
\end{Verbatim}

\begin{Verbatim}[fontsize=\footnotesize,breaklines=true,breakanywhere=true]
Why the wrong ranking is wrong:
This ranking is wrong because it promotes unsupported alternatives, includes overly specific labels without evidence from the description, and fails to place the description-supported type first.
\end{Verbatim}

\begin{Verbatim}[fontsize=\footnotesize,breaklines=true,breakanywhere=true]
Now classify:
\end{Verbatim}

\begin{Verbatim}[fontsize=\footnotesize,breaklines=true,breakanywhere=true]
Entity: framboid
Description: minute spheroidal grain of pyrite or an aggregate of such grains, commonly formed at low temperatures in oxygen-poor mud sediments containing organic matter, especially in marine mud deposits
\end{Verbatim}

\begin{Verbatim}[fontsize=\footnotesize,breaklines=true,breakanywhere=true]
Answer:
\end{Verbatim}

\clearpage

\section{Contrastive graph}

\begin{Verbatim}[fontsize=\footnotesize,breaklines=true,breakanywhere=true]
List valid entity-type labels, most specific first, then broader valid types. Use concise labels, not entity names or descriptions. Include only evidence-supported types; do not repeat labels or pad the list.
\end{Verbatim}

\begin{Verbatim}[fontsize=\footnotesize,breaklines=true,breakanywhere=true]
The context surrounds an omitted target type; it is evidence, not answer options. Parents are broader, children narrower, and siblings share a parent. Infer the type using the description and these relations. Copy a neighbour only if valid for this entity. Siblings need not be disjoint; predict a child only when the description independently supports it.
\end{Verbatim}

\begin{Verbatim}[fontsize=\footnotesize,breaklines=true,breakanywhere=true]
Briefly use parents to explain scope, compare sibling alternatives, and check children for unsupported specificity. Justify the final specific-to-broad order using the description and taxonomy evidence.
\end{Verbatim}

\begin{Verbatim}[fontsize=\footnotesize,breaklines=true,breakanywhere=true]
Contrast the inferred type with sibling and child alternatives. Reject an alternative only when the evidence warrants it.
\end{Verbatim}

\begin{Verbatim}[fontsize=\footnotesize,breaklines=true,breakanywhere=true]
Return JSON only in exactly this form:
{"explanation": "contrastive taxonomy-grounded explanation", "types": ["type 1", "type 2", "..."]}
\end{Verbatim}

\begin{Verbatim}[fontsize=\footnotesize,breaklines=true,breakanywhere=true]
Example:
\end{Verbatim}

\begin{Verbatim}[fontsize=\footnotesize,breaklines=true,breakanywhere=true]
Entity: Rod Fergusson
Description: Canadian video game producer
\end{Verbatim}

\begin{Verbatim}[fontsize=\footnotesize,breaklines=true,breakanywhere=true]
Taxonomy context:
- Direct parent classes: producer
- Sibling classes: creative producer, co-producer, line producer
- Direct child classes: None
\end{Verbatim}

\begin{Verbatim}[fontsize=\footnotesize,breaklines=true,breakanywhere=true]
Correct response:
{"explanation": "The description \"Canadian video game producer\" supports video game producer as the most specific type. In the taxonomy, its direct parent is producer, which constrains the broader semantic category. The sibling classes are nearby alternatives, but classes such as creative producer, co-producer, line producer are not as well supported by the description. The listed child context is None; a child should be selected only if the description independently supports that narrower meaning. The remaining returned types are broader valid classes ranked in decreasing specificity.", "types": ["video game producer", "producer", "media professional", "professional", "worker", "person", "person or organization", "corporate body", "object", "entity"]}
\end{Verbatim}

\begin{Verbatim}[fontsize=\footnotesize,breaklines=true,breakanywhere=true]
Wrong ranking:
{"types": ["producer", "creative producer", "co-producer", "line producer", "radio producer", "DJ producer", "professional", "worker", "person", "entity"]}
\end{Verbatim}

\begin{Verbatim}[fontsize=\footnotesize,breaklines=true,breakanywhere=true]
Why the wrong ranking is wrong:
This ranking is wrong because it treats taxonomy neighbours as answer options, promotes classes that the entity description does not support, and omits the more specific description-supported target type.
\end{Verbatim}

\begin{Verbatim}[fontsize=\footnotesize,breaklines=true,breakanywhere=true]
Now classify:
\end{Verbatim}

\begin{Verbatim}[fontsize=\footnotesize,breaklines=true,breakanywhere=true]
Entity: framboid
Description: minute spheroidal grain of pyrite or an aggregate of such grains, commonly formed at low temperatures in oxygen-poor mud sediments containing organic matter, especially in marine mud deposits
\end{Verbatim}

\begin{Verbatim}[fontsize=\footnotesize,breaklines=true,breakanywhere=true]
Taxonomy context:
- Direct parent classes: chemical substance
- Sibling classes: chemical product, pesticide, aroma compound
- Direct child classes: None
\end{Verbatim}

\begin{Verbatim}[fontsize=\footnotesize,breaklines=true,breakanywhere=true]
Answer:
\end{Verbatim}

\clearpage

\section{Parents only}

\begin{Verbatim}[fontsize=\footnotesize,breaklines=true,breakanywhere=true]
List valid entity-type labels, most specific first, then broader valid types. Use concise labels, not entity names or descriptions. Include only evidence-supported types; do not repeat labels or pad the list.
\end{Verbatim}

\begin{Verbatim}[fontsize=\footnotesize,breaklines=true,breakanywhere=true]
The taxonomy context contains only direct parent classes of an omitted target type. Parent classes are broader than the omitted target and can constrain its semantic scope. They are structural clues, not answer options. Do not copy a parent into the answer unless it is independently a valid type for the entity.
\end{Verbatim}

\begin{Verbatim}[fontsize=\footnotesize,breaklines=true,breakanywhere=true]
Return JSON only in exactly this form:
{"types": ["type 1", "type 2", "..."]}
\end{Verbatim}

\begin{Verbatim}[fontsize=\footnotesize,breaklines=true,breakanywhere=true]
Example:
\end{Verbatim}

\begin{Verbatim}[fontsize=\footnotesize,breaklines=true,breakanywhere=true]
Entity: Rod Fergusson
Description: Canadian video game producer
\end{Verbatim}

\begin{Verbatim}[fontsize=\footnotesize,breaklines=true,breakanywhere=true]
Taxonomy context:
- Direct parent classes: producer
\end{Verbatim}

\begin{Verbatim}[fontsize=\footnotesize,breaklines=true,breakanywhere=true]
Answer:
{"types": ["video game producer", "producer", "media professional", "professional", "worker", "person", "person or organization", "corporate body", "object", "entity"]}
\end{Verbatim}

\begin{Verbatim}[fontsize=\footnotesize,breaklines=true,breakanywhere=true]
Now classify:
\end{Verbatim}

\begin{Verbatim}[fontsize=\footnotesize,breaklines=true,breakanywhere=true]
Entity: framboid
Description: minute spheroidal grain of pyrite or an aggregate of such grains, commonly formed at low temperatures in oxygen-poor mud sediments containing organic matter, especially in marine mud deposits
\end{Verbatim}

\begin{Verbatim}[fontsize=\footnotesize,breaklines=true,breakanywhere=true]
Taxonomy context:
- Direct parent classes: chemical substance
\end{Verbatim}

\begin{Verbatim}[fontsize=\footnotesize,breaklines=true,breakanywhere=true]
Answer:
\end{Verbatim}

\clearpage

\section{Siblings only}

\begin{Verbatim}[fontsize=\footnotesize,breaklines=true,breakanywhere=true]
List valid entity-type labels, most specific first, then broader valid types. Use concise labels, not entity names or descriptions. Include only evidence-supported types; do not repeat labels or pad the list.
\end{Verbatim}

\begin{Verbatim}[fontsize=\footnotesize,breaklines=true,breakanywhere=true]
The taxonomy context contains only sibling classes of an omitted target type. Siblings share a direct parent with the omitted target and are nearby alternatives, but they are not necessarily logically disjoint. They are structural clues, not answer options. Do not copy a sibling into the answer unless it is independently a valid type for the entity.
\end{Verbatim}

\begin{Verbatim}[fontsize=\footnotesize,breaklines=true,breakanywhere=true]
Return JSON only in exactly this form:
{"types": ["type 1", "type 2", "..."]}
\end{Verbatim}

\begin{Verbatim}[fontsize=\footnotesize,breaklines=true,breakanywhere=true]
Example:
\end{Verbatim}

\begin{Verbatim}[fontsize=\footnotesize,breaklines=true,breakanywhere=true]
Entity: Rod Fergusson
Description: Canadian video game producer
\end{Verbatim}

\begin{Verbatim}[fontsize=\footnotesize,breaklines=true,breakanywhere=true]
Taxonomy context:
- Sibling classes: creative producer, co-producer, line producer
\end{Verbatim}

\begin{Verbatim}[fontsize=\footnotesize,breaklines=true,breakanywhere=true]
Answer:
{"types": ["video game producer", "producer", "media professional", "professional", "worker", "person", "person or organization", "corporate body", "object", "entity"]}
\end{Verbatim}

\begin{Verbatim}[fontsize=\footnotesize,breaklines=true,breakanywhere=true]
Now classify:
\end{Verbatim}

\begin{Verbatim}[fontsize=\footnotesize,breaklines=true,breakanywhere=true]
Entity: framboid
Description: minute spheroidal grain of pyrite or an aggregate of such grains, commonly formed at low temperatures in oxygen-poor mud sediments containing organic matter, especially in marine mud deposits
\end{Verbatim}

\begin{Verbatim}[fontsize=\footnotesize,breaklines=true,breakanywhere=true]
Taxonomy context:
- Sibling classes: chemical product, pesticide, aroma compound
\end{Verbatim}

\begin{Verbatim}[fontsize=\footnotesize,breaklines=true,breakanywhere=true]
Answer:
\end{Verbatim}

\clearpage

\section{Children only}

\begin{Verbatim}[fontsize=\footnotesize,breaklines=true,breakanywhere=true]
List valid entity-type labels, most specific first, then broader valid types. Use concise labels, not entity names or descriptions. Include only evidence-supported types; do not repeat labels or pad the list.
\end{Verbatim}

\begin{Verbatim}[fontsize=\footnotesize,breaklines=true,breakanywhere=true]
The taxonomy context contains only direct child classes of an omitted target type. Child classes are narrower than the omitted target and can help indicate its specificity boundary. They are structural clues, not answer options. Do not predict a child merely because it is listed; use a narrower child only when the entity description independently supports it.
\end{Verbatim}

\begin{Verbatim}[fontsize=\footnotesize,breaklines=true,breakanywhere=true]
Return JSON only in exactly this form:
{"types": ["type 1", "type 2", "..."]}
\end{Verbatim}

\begin{Verbatim}[fontsize=\footnotesize,breaklines=true,breakanywhere=true]
Example:
\end{Verbatim}

\begin{Verbatim}[fontsize=\footnotesize,breaklines=true,breakanywhere=true]
Entity: Rod Fergusson
Description: Canadian video game producer
\end{Verbatim}

\begin{Verbatim}[fontsize=\footnotesize,breaklines=true,breakanywhere=true]
Taxonomy context:
- Direct child classes: None
\end{Verbatim}

\begin{Verbatim}[fontsize=\footnotesize,breaklines=true,breakanywhere=true]
Answer:
{"types": ["video game producer", "producer", "media professional", "professional", "worker", "person", "person or organization", "corporate body", "object", "entity"]}
\end{Verbatim}

\begin{Verbatim}[fontsize=\footnotesize,breaklines=true,breakanywhere=true]
Now classify:
\end{Verbatim}

\begin{Verbatim}[fontsize=\footnotesize,breaklines=true,breakanywhere=true]
Entity: framboid
Description: minute spheroidal grain of pyrite or an aggregate of such grains, commonly formed at low temperatures in oxygen-poor mud sediments containing organic matter, especially in marine mud deposits
\end{Verbatim}

\begin{Verbatim}[fontsize=\footnotesize,breaklines=true,breakanywhere=true]
Taxonomy context:
- Direct child classes: None
\end{Verbatim}

\begin{Verbatim}[fontsize=\footnotesize,breaklines=true,breakanywhere=true]
Answer:
\end{Verbatim}
